% Part: lambda-calculus
% Chapter: introduction
% Section: minimization

\documentclass[../../../include/open-logic-section]{subfiles}

\begin{document}

\olfileid{lam}{int}{min}
\olsection{الدوال \usetoken{S}{lambda definable} مغلقة تحت التصغير}

\begin{lem}
لتكن~$f(x,y)$ !!{lambda definable}، ولتعرّف~$g$ بالمعادلة
\[
g(x) \simeq \umin{y}{f(x,y)}.
\]
فإن~$g$ أيضًا !!{lambda definable}.
\end{lem}

\begin{proof}
نجمل وجه البرهان أولًا. إذا أعطينا~$x$، استعنا بحد النقطة الثابتة
اللامبداوي~$Y$ لتعريف دالة~$h_x(n)$ تطلب~$y$ ابتداءً من~$n$.
وعندئذ يكون~$g(x)$ هو~$h_x(0)$. أما~$h_x$ فنطلبها حلًا لمعادلة النقطة الثابتة الآتية:
\[
h_x(n) \simeq
\begin{cases}
n & \text{إذا كان $f(x,n) = 0$} \\
h_x(n+1) & \text{في غير ذلك.}
\end{cases}
\]

وأما التفصيل، فإن~$f$، لكونها قابلة للتعريف بلامبدا، هي
!!{lambda defined} بحد~$F$. وقد سبق أن لدينا حد لامبدا~$D$ يحقق
$D(M, N, \bar{0}) \red M$ و$D(M, N, \bar{1}) \red N$. نثبت~$x$ في هذه الخطوة؛ ولكي !!{lambda define} حدٌ الدالة~$h_x$،
نطلب حدًا~$H$ يعتمد على~$x$ ويحقق
\[
H(\num{n}) \equiv D(\num{n}, H(S(\num{n})), F(x, \num{n})).
\]
ويتيح لنا حد النقطة الثابتة~$Y$ هذا الإنشاء. نأخذ أولًا~$U$ مساويًا للحد
\[
\lambd[h][\lambd[z][D(z,(h(Sz)),F(x,z))]],
\]
ثم نأخذ~$H$ مساويًا للحد~$YU$. والمتغير الحر الوحيد في~$H$ هو~$x$.
فلنثبت أن~$H$ يفي بالمعادلة المذكورة.

من تعريف~$Y$ ينتج
\[
H = YU \equiv U(YU) = U(H).
\]
ومن ثم، لكل عدد طبيعي~$n$،
\begin{align*}
H(\num{n}) & \equiv U(H, \num{n}) \\
& \red D(\num{n}, H(S(\num{n})), F(x, \num{n})),
\end{align*}
وهذا هو المطلوب. ثم إذا وضعنا العدد~$\num{m}$ موضع~$x$ في السطر الأخير،
فإن التعبير يختزل إلى~$\num{n}$ متى اختزل~$F(\num{m}, \num{n})$ إلى~$\num{0}$،
وإلى~$H(S(\num{n}))$ متى اختزل~$F(\num{m}, \num{n})$ إلى عدد آخر.

ويبقى أن نأخذ~$G$ مساويًا لـ$\lambd[x][H(\num 0)]$. فهذا الحد
$G$ !!{lambda define}s~$g$: فإن~$G(\num m)$، عند كل~$m$،
يختزل إلى~$\overline {g(m)}$ إن كانت~$g(m)$ معرّفة، وأما إن لم تكن معرّفة فلا صورة سوية له.
\end{proof}

\end{document}

